\expandafter\def\csname HistoricalNoteText3\endcsname{Der analytischen Vollständigkeit wegen bemerke ich, dass in diesem Falle der Himmelskörper in
$\left\{\begin{array}{c}\AscendingNode\\\DescendingNode\end{array}\right\}$ ist, je nachdem $\sin(\odot-\alpha)$ und $\sin(\alpha-\AscendingNode)$
$\left\{\begin{array}{l}\text{einerlei}\\\text{entgegengesetzte}\end{array}\right\}$ Zeichen haben.}
\clearpage
\begingroup
\pagestyle{empty}
\renewcommand{\footnoterule}{\kern-3pt\hrule width .285\linewidth\kern2.6pt}
\renewcommand{\footnotesize}{\fontsize{8.4}{10.6}\selectfont}
\newcommand{\GaussVIGathering}[2]{%
  \fancypagestyle{gaussVIGathering}{%
    \fancyhf{}%
    \renewcommand{\headrulewidth}{0pt}%
    \renewcommand{\footrulewidth}{0pt}%
    \fancyfoot[C]{\makebox[\textwidth][l]{%
      \hspace*{.025\textwidth}{\footnotesize\bfseries #1}\hfill
      {\normalsize\bfseries #2}\hspace*{.18\textwidth}}}%
  }%
  \thispagestyle{gaussVIGathering}%
}
\SourceLeaf{92}%
\pdfbookmark[1]{Vorschriften zur Bestimmung heliocentrischer Orte}{vorschriften1802}
\vspace*{.045\textheight}
{\centering\fontsize{7.85}{10}\selectfont
VORSCHRIFTEN, UM AUS DER GEOCENTRISCHEN LÄNGE UND BREITE EINES\par
HIMMELSKÖRPERS, DEM ORTE SEINES KNOTENS, DER NEIGUNG DER BAHN, DER LÄNGE\par
DER SONNE UND IHREM ABSTANDE VON DER ERDE ABZULEITEN:\par
DES HIMMELSKÖRPERS HELIOCENTRISCHE LÄNGE IN DER BAHN, WAHREN ABSTAND\par
VON DER SONNE UND WAHREN ABSTAND VON DER ERDE.\par
}
\bigskip
\noindent\makebox[\linewidth][c]{%
  \parbox{.82\linewidth}{%
    \centering\fontsize{8.7}{11}\selectfont
    \hrule height .3pt\smallskip
    Monatliche Correspondenz zur Beförderung der Erd- und Himmels-Kunde,\par
    herausgegeben vom Freiherrn \textsc{von Zach}. Juni 1802.\par
    \smallskip\hrule height .3pt
  }%
}\par
\medskip
{\centering\textit{Bedeutung der Zeichen:}\par}
\noindent\begin{tabular}{@{}c@{\hspace{.025\linewidth}}|@{\hspace{.025\linewidth}}c@{}}
\begin{minipage}[t]{.474\linewidth}\small
\begin{center}\textit{Gegeben;}\end{center}
\begin{tabular}{@{}lp{.78\linewidth}@{}}
$\AscendingNode$ & \raggedright Länge des aufsteigenden Knotens.\tabularnewline
$\odot$ & \raggedright Länge der Sonne.\tabularnewline
$\alpha$ & \raggedright Geocentrische Länge des Himmelskörpers.\tabularnewline
$\ObservedBeta$ & \raggedright Geocentrische Breite.\tabularnewline
$i$ & \raggedright Neigung der Bahn.\tabularnewline
$R$ & \raggedright Abstand der Sonne von der Erde.
\end{tabular}
\end{minipage}&
\begin{minipage}[t]{.474\linewidth}\small
\begin{center}\textit{Gesucht:}\end{center}
\begin{tabular}{@{}lp{.78\linewidth}@{}}
$\upsilon$ & \raggedright heliocentr. Länge des Himmelskörpers in der Bahn.\tabularnewline
$r$ & \raggedright Wahrer Abstand von der Sonne.\tabularnewline
$\Delta$ & \raggedright Wahrer Abstand von der Erde.
\end{tabular}
\[
\left.\begin{array}{l}A\\B\\C\ \text{etc.}\end{array}\right\}\text{ Hülfswinkel.}
\]
\end{minipage}
\end{tabular}\par

\smallskip
{\centering I.\par}
{\fontsize{9}{11}\selectfont
\setlength{\abovedisplayskip}{6pt}
\setlength{\belowdisplayskip}{6pt}
\setlength{\abovedisplayshortskip}{6pt}
\setlength{\belowdisplayshortskip}{6pt}
\[
\begin{array}{@{}l@{\quad}r@{\;}l@{\qquad}r@{\;}l@{}}
1^{0} & \dfrac{\cos(\odot-\AscendingNode)\operatorname{tang}\ObservedBeta}{\sin(\odot-\alpha)} & {}=\operatorname{tang}A & \dfrac{\sin A\operatorname{tang}(\odot-\AscendingNode)}{\sin(A+i)} & {}=\operatorname{tang}(\upsilon-\AscendingNode)\\[4pt]
2^{0} & \dfrac{\sin(\odot-\alpha)\operatorname{tang}i}{\cos(\odot-\AscendingNode)} & {}=\operatorname{tang}B & \dfrac{\cos B\sin\ObservedBeta\operatorname{tang}(\odot-\AscendingNode)}{\sin(B+\ObservedBeta)\cos i} & {}=\operatorname{tang}(\upsilon-\AscendingNode)\\[4pt]
3^{0} & \dfrac{\sin(\odot-\AscendingNode)\operatorname{tang}\ObservedBeta}{\sin(\odot-\alpha)\operatorname{tang}i} & {}=\operatorname{tang}C & \dfrac{\sin C\sin(\odot-\AscendingNode)}{\sin(C+\odot-\AscendingNode)\cos i} & {}=\operatorname{tang}(\upsilon-\AscendingNode)\\[4pt]
4^{0} & \dfrac{\cos(\odot-\AscendingNode)\operatorname{tang}\ObservedBeta}{\cos(\odot-\alpha)\operatorname{tang}i} & {}=\operatorname{tang}D & \dfrac{\sin D\operatorname{tang}(\odot-\AscendingNode)\cos(\odot-\alpha)}{\sin(D+\odot-\alpha)\cos i} & {}=\operatorname{tang}(\upsilon-\AscendingNode)
\end{array}
\]
}

\textit{Anmerkung:} Da Winkel, die um $180^{0}$ verschieden sind, einerlei Tangenten
haben, so ist hier noch eine Vorschrift nöthig, wie die durch ihre Tangenten
bestimmten Winkel $A$, $B$, $C$ etc. und $\upsilon-\AscendingNode$ angesetzt werden müssen. Den
Winkel $\upsilon-\AscendingNode$ hat man allezeit zwischen 0 und $180^{0}$ anzunehmen, wenn $\ObservedBeta$
\par\clearpage
\noindent\SourceLeaf{93}%
\makebox[12mm][l]{{\large\bfseries 88}}%
\parbox[t]{\dimexpr\linewidth-24mm\relax}{\centering\fontsize{7.05}{9.2}\selectfont\bfseries VORSCHRIFTEN UM AUS DER GEOCENTRISCHEN LÄNGE UND BREITE\par}%
\makebox[12mm][r]{}\par\nobreak\smallskip
\noindent positiv (nördlich) ist; ist hingegen die Breite südlich, so muss $\upsilon-\AscendingNode$ zwischen
$180^{0}$ und $360^{0}$, oder was einerlei ist, zwischen $-180^{0}$ und 0 fallen. Ist
$\ObservedBeta=0$, so ist der Himmelskörper in einem Knoten, und man wird nie zweifelhaft
sein, ob es $\AscendingNode$ oder $\DescendingNode$ ist\HistoricalNote{3}. Die Hülfswinkel, $A$, $B$, $C$, $D$ aber, so wie
die folgenden $E$, $F$ etc., kann man in dieser Hinsicht ganz nach Belieben ansetzen;
wobei es sich jedoch von selbst versteht, dass man auf die Zeichen $\pm$
gehörige Rücksicht nehme; ich habe sie in folgendem Beispiele immer zwischen
$-90^{0}$ und $+90^{0}$ genommen. — Logarithmen, deren zugehörige Grössen negativ
sind, habe ich durch ein beigeschriebenes n ausgezeichnet.

{\centering II.\par}
{\fontsize{9}{11}\selectfont
\setlength{\abovedisplayskip}{6pt}
\setlength{\belowdisplayskip}{6pt}
\setlength{\abovedisplayshortskip}{6pt}
\setlength{\belowdisplayshortskip}{6pt}
\[
\begin{array}{@{}l@{\quad}r@{\;}l@{\qquad}r@{\;}l@{}}
5^{0} & \dfrac{\operatorname{tang}\ObservedBeta}{\sin(\alpha-\AscendingNode)} & {}=\operatorname{tang}E & \dfrac{\sin E\sin(\odot-\AscendingNode)}{\sin(i-E)\sin(\upsilon-\AscendingNode)} & {}=\dfrac{r}{R}\\[4pt]
6^{0} & \operatorname{tang}i\sin(\alpha-\AscendingNode) & {}=\operatorname{tang}F & \dfrac{\cos F\sin(\odot-\AscendingNode)\sin\ObservedBeta}{\sin(F-\ObservedBeta)\sin(\upsilon-\AscendingNode)\cos i} & {}=\dfrac{r}{R}\\[4pt]
7^{0} & \cos i\operatorname{tang}(\upsilon-\AscendingNode) & {}=\operatorname{tang}G & \dfrac{\cos G\sin(\odot-\alpha)}{\sin(\alpha-\AscendingNode-G)\cos(\upsilon-\AscendingNode)} & {}=\dfrac{r}{R}\\[4pt]
8^{0} & \dfrac{\operatorname{tang}(\alpha-\AscendingNode)}{\cos i} & {}=\operatorname{tang}H & \dfrac{\sin H\sin(\odot-\alpha)}{\sin(H-(\upsilon-\AscendingNode))\sin(\alpha-\AscendingNode)} & {}=\dfrac{r}{R}\\[4pt]
9^{0} & \dfrac{\operatorname{tang}\ObservedBeta}{\sin i\cos(\alpha-\AscendingNode)} & {}=\operatorname{tang}I & \dfrac{\sin I\cos(\odot-\AscendingNode)}{\sin(\upsilon-\AscendingNode-I)} & {}=\dfrac{r}{R}\\[4pt]
10^{0} & \sin i\cos(\alpha-\AscendingNode)\operatorname{tang}(\upsilon-\AscendingNode) & {}=\operatorname{tang}K & \dfrac{\cos K\sin\ObservedBeta\cos(\odot-\AscendingNode)}{\sin(K-\ObservedBeta)\cos(\upsilon-\AscendingNode)} & {}=\dfrac{r}{R}\\[4pt]
11^{0} & \dfrac{\sin C\sin(\odot-\alpha)}{\cos(C+\odot-\alpha)\operatorname{tang}(\odot-\AscendingNode)\cos i} & {}=\operatorname{tang}L & \dfrac{\sin L}{\sin(\upsilon-\AscendingNode-L)\cos(\odot-\AscendingNode)} & {}=\dfrac{r}{R}\\[4pt]
12^{0} & \dfrac{\sin D\cos(\odot-\AscendingNode)}{\cos(D+\odot-\AscendingNode)\cos i} & {}=\operatorname{tang}M & \dfrac{\sin M}{\sin(\upsilon-\AscendingNode-M)\cos(\odot-\AscendingNode)} & {}=\dfrac{r}{R}
\end{array}
\]
}
{\centering III.\par}
{\fontsize{9}{11}\selectfont
\setlength{\abovedisplayskip}{6pt}
\setlength{\belowdisplayskip}{6pt}
\setlength{\abovedisplayshortskip}{6pt}
\setlength{\belowdisplayshortskip}{6pt}
\[
\begin{array}{@{}l@{\quad}l@{}}
13^{0} & \dfrac{r\sin(\upsilon-\AscendingNode)\sin i}{\sin\ObservedBeta}=\Delta\\[4pt]
14^{0} & \dfrac{R\sin E\sin(\odot-\AscendingNode)\sin i}{\sin(i-E)\sin\ObservedBeta}=\dfrac{R\cos E\sin(\odot-\AscendingNode)\sin i}{\sin(i-E)\sin(\alpha-\AscendingNode)\cos\ObservedBeta}=\Delta\\[4pt]
15^{0} & \dfrac{R\cos F\sin(\odot-\AscendingNode)\operatorname{tang}i}{\sin(F-\ObservedBeta)}=\dfrac{R\sin F\sin(\odot-\AscendingNode)\sin(\alpha-\AscendingNode)}{\sin(F-\ObservedBeta)}=\Delta
\end{array}
\]
}

Und so lassen sich noch mehrere Ausdrücke für $\Delta$ aus der Verbindung
von $13^{0}$ mit allen Formeln II ableiten.

\par\clearpage
\noindent\SourceLeaf{94}%
\parbox[t]{\dimexpr\linewidth-12mm\relax}{\centering\fontsize{6.2}{8.8}\selectfont\bfseries EINES HIMMELSKÖRPERS DEM ORTE SEINES KNOTENS DER NEIGUNG DER BAHN ETC.\par}%
\makebox[12mm][r]{{\large\bfseries 89}}\par\nobreak\smallskip
{\centering\textit{Beispiel.}\par}
\[
\begin{array}{rlrrr}
\AscendingNode&=&80^{0}&59'&12''07\\
\odot&=&281&1&34.99\\
\alpha&=&53&23&2.46\\
i&=&10&37&9.55
\end{array}
\]
\[
\begin{array}{@{}r@{\;}l@{\quad}l@{}}
\log\operatorname{tang}\ObservedBeta & {}=8.7349698\,\mathrm n & \ObservedBeta=-3^{0}\ 6'\ 33''561\quad\text{negativ oder südlich}\\
\log R & {}=9.9926158 &
\end{array}
\]
\[
\begin{array}{rlrrr}
\text{Folglich }\odot-\AscendingNode&=&200^{0}&2'&22''92\\
\odot-\alpha&=&227&38&32.53\\
\alpha-\AscendingNode&=&-27&36&9.61
\end{array}
\]

\par\medskip\noindent\begin{minipage}[t]{\textwidth}
{\centering $1^{0}$.\par}
\smallskip\fontsize{8.7}{10.5}\selectfont
\noindent\begin{tabular*}{\linewidth}{@{}l@{\enspace}l@{\extracolsep{\fill}\quad}l@{\extracolsep{0pt}\enspace}r@{\enspace}r@{\enspace}r@{\enspace}r@{}}
$\log\operatorname{tang}\ObservedBeta\enspace .\ .\ .\ .\ .$ & $8.7349698\,\mathrm n$ & & & & &\\
$\log\cos(\odot-\AscendingNode)\enspace .\ .\ .$ & $9.9728762\,\mathrm n$ & & & & &\\
$\text{Compl.}\ \log\sin(\odot-\alpha)\enspace .$ & $0.1313827\,\mathrm n$ & & & & &\\
\cline{1-2}
$\log\operatorname{tang}A\enspace .\ .\ .\ .\ .$ & $8.8392287\,\mathrm n$ & folglich & $A=-$ & $3^{0}$ & $57'$ & $2''136$\\
\cline{1-2}
$\log\sin A\enspace .\ .\ .\ .\ .$ & $8.8381955\,\mathrm n$ &  & $A+i=$ & $6$ & $40$ & $7.414$\\
$\log\operatorname{tang}(\odot-\AscendingNode)\enspace .\ .\ .$ & $9.5620014$ & & & & &\\
$\text{Compl.}\ \log\sin(A+i)\enspace .$ & $0.9350608$ & & & & &\\
\cline{1-2}
$\log\operatorname{tang}(\upsilon-\AscendingNode)\enspace .\ .\ .$ & $9.3352577\,\mathrm n$ & ferner & $\upsilon-\AscendingNode=-$ & $12^{0}$ & $12'$ & $37''942$\\
 &  & also & $\upsilon=.$ & $68$ & $46$ & $34.128$\\
\cline{3-7}
\end{tabular*}
\end{minipage}\par
\par\medskip\noindent\begin{minipage}[t]{\textwidth}
{\centering $2^{0}$.\par}
\smallskip\fontsize{8.7}{10.5}\selectfont
\noindent\begin{tabular*}{\linewidth}{@{}l@{\enspace}l@{\extracolsep{\fill}\quad}l@{\extracolsep{0pt}\enspace}r@{\enspace}r@{\enspace}r@{\enspace}r@{}}
$\log\sin(\odot-\alpha)\enspace .\ .\ .\ .$ & $9.8686173\,\mathrm n$ & & & & &\\
$\log\operatorname{tang}i\enspace .\ .\ .\ .\ .$ & $9.2729872$ & & & & &\\
$\text{Compl.}\ \log\cos(\odot-\AscendingNode)\enspace .$ & $0.0271238\,\mathrm n$ & & & & &\\
\cline{1-2}
$\log\operatorname{tang}B\enspace .\ .\ .\ .\ .$ & $9.1687283$ & folglich & $B=$ & $8^{0}$ & $23'$ & $21''888$\\
\cline{1-2}
$\log\cos B\enspace .\ .\ .\ .\ .$ & $9.9953277$ &  & $B+\ObservedBeta=$ & $5$ & $16$ & $48.327$\\
$\log\sin\ObservedBeta\enspace .\ .\ .\ .\ .\ .$ & $8.7343300\,\mathrm n$ & & & & &\\
$\log\operatorname{tang}(\odot-\AscendingNode)\enspace .\ .\ .$ & $9.5620014$ & & & & &\\
$\text{Compl.}\ \log\sin(B+\ObservedBeta)\enspace .$ & $1.0360961$ & & & & &\\
$\text{Compl.}\ \log\cos i\enspace .\ .\ .$ & $0.0075025$ & & & & &\\
\cline{1-2}
$\log\operatorname{tang}(\upsilon-\AscendingNode)\enspace .\ .\ .$ & $9.3352577\,\mathrm n$ & \multicolumn{5}{@{\quad}l@{}}{wie oben}\\
\end{tabular*}
\end{minipage}\par
\par\GaussVIGathering{VI.}{12}
\par\clearpage
\noindent\SourceLeaf{95}%
\makebox[12mm][l]{{\large\bfseries 90}}%
\parbox[t]{\dimexpr\linewidth-24mm\relax}{\centering\fontsize{7.05}{9.2}\selectfont\bfseries VORSCHRIFTEN UM AUS DER GEOGRAPHISCHEN LÄNGE UND BREITE\par}%
\makebox[12mm][r]{}\par\nobreak\smallskip
\par\medskip\noindent\begin{minipage}[t]{\textwidth}
{\centering $3^{0}$.\par}
\smallskip\fontsize{8.7}{10.5}\selectfont
\noindent\begin{tabular*}{\linewidth}{@{}l@{\enspace}l@{\extracolsep{\fill}\quad}l@{\extracolsep{0pt}\enspace}r@{\enspace}r@{\enspace}r@{\enspace}r@{}}
$\log\sin(\odot-\AscendingNode)\enspace .\ .\ .$ & $9.5348776\,\mathrm n$ & & & & &\\
$\log\operatorname{tang}\ObservedBeta\enspace .\ .\ .\ .\ .$ & $8.7349698\,\mathrm n$ & & & & &\\
$\text{Compl.}\ \log\sin(\odot-\alpha)\enspace .$ & $0.1313827\,\mathrm n$ & & & & &\\
$\text{Compl.}\ \log\operatorname{tang}i\enspace .\ .\ .$ & $0.7270128$ & & & & &\\
\cline{1-2}
$\log\operatorname{tang}C\enspace .\ .\ .\ .\ .$ & $9.1282429\,\mathrm n$ & also & $C=-$ & $7^{0}$ & $39'$ & $7''056$\\
\cline{1-2}
$\log\sin C\enspace .\ .\ .\ .\ .$ & $9.1243583\,\mathrm n$ &  & $C+\odot-\AscendingNode=$ & $192$ & $23$ & $15.864$\\
$\log\sin(\odot-\AscendingNode)\enspace .\ .\ .$ & $9.5348776\,\mathrm n$ & & & & &\\
$\text{Compl.}\ \log\sin(C+\odot-\AscendingNode)$ & $0.6685194\,\mathrm n$ & & & & &\\
$\text{Compl.}\ \log\cos i\enspace .\ .\ .$ & $0.0075025$ & & & & &\\
\cline{1-2}
$\log\operatorname{tang}(\upsilon-\AscendingNode)\enspace .\ .\ .$ & $9.3352578\,\mathrm n$ & \multicolumn{5}{@{\quad}l@{}}{wie vorhin.}\\
\end{tabular*}
\end{minipage}\par
\par\medskip\noindent\begin{minipage}[t]{\textwidth}
{\centering $4^{0}$.\par}
\smallskip\fontsize{8.7}{10.5}\selectfont
\noindent\begin{tabular*}{\linewidth}{@{}l@{\enspace}l@{\extracolsep{\fill}\quad}l@{\extracolsep{0pt}\enspace}r@{\enspace}r@{\enspace}r@{\enspace}r@{}}
$\log\cos(\odot-\AscendingNode)\enspace .\ .\ .$ & $9.9728762\,\mathrm n$ & & & & &\\
$\log\operatorname{tang}\ObservedBeta\enspace .\ .\ .\ .\ .$ & $8.7349698\,\mathrm n$ & & & & &\\
$\text{Compl.}\ \log\cos(\odot-\alpha)\enspace .$ & $0.1714973\,\mathrm n$ & & & & &\\
$\text{Compl.}\ \log\operatorname{tang}i\enspace .\ .\ .$ & $0.7270128$ & & & & &\\
\cline{1-2}
$\log\operatorname{tang}D\enspace .\ .\ .\ .\ .$ & $9.6063561\,\mathrm n$ & also & $D=-$ & $21^{0}$ & $59'$ & $51''182$\\
\cline{1-2}
$\log\sin D\enspace .\ .\ .\ .\ .$ & $9.5735295\,\mathrm n$ &  & $D+\odot-\alpha=$ & $205$ & $38$ & $41.348$\\
$\log\operatorname{tang}(\odot-\AscendingNode)\enspace .\ .\ .$ & $9.5620014$ & & & & &\\
$\log\cos(\odot-\alpha)\enspace .\ .\ .\ .$ & $9.8285027\,\mathrm n$ & & & & &\\
$\text{Compl.}\ \log\sin(D+\odot-\alpha)$ & $0.3637217\,\mathrm n$ & & & & &\\
$\text{Compl.}\ \log\cos i\enspace .\ .\ .$ & $0.0075025$ & & & & &\\
\cline{1-2}
$\log\operatorname{tang}(\upsilon-\AscendingNode)\enspace .\ .\ .\ .$ & $9.3352578\,\mathrm n$ & \multicolumn{5}{@{\quad}l@{}}{wie oben.}\\
\end{tabular*}
\end{minipage}\par
\par\medskip\noindent\begin{minipage}[t]{\textwidth}
{\centering $5^{0}$.\par}
\smallskip\fontsize{8.7}{10.5}\selectfont
\noindent\begin{tabular*}{\linewidth}{@{}l@{\enspace}l@{\extracolsep{\fill}\quad}l@{\extracolsep{0pt}\enspace}r@{\enspace}r@{\enspace}r@{\enspace}r@{}}
$\log\operatorname{tang}\ObservedBeta\enspace .\ .\ .\ .\ .$ & $8.7349698\,\mathrm n$ & & & & &\\
$\log\sin(\alpha-\AscendingNode)\enspace .\ .\ .\ .$ & $9.6658973\,\mathrm n$ & & & & &\\
\cline{1-2}
$\log\operatorname{tang}E\enspace .\ .\ .\ .\ .$ & $9.0690725$ & also & $E=$ & $6^{0}$ & $41'$ & $12''412$\\
\cline{1-2}
$\log\sin E\enspace .\ .\ .\ .\ .$ & $9.0661081$ &  & $i-E=$ & $3$ & $55$ & $57.138$\\
$\log\sin(\odot-\AscendingNode)\enspace .\ .\ .\ .$ & $9.5348776\,\mathrm n$ & & & & &\\
$\text{Compl.}\ \log\sin(i-E)\enspace .$ & $1.1637907$ & & & & &\\
$\text{Compl.}\ \log\sin(\upsilon-\AscendingNode)\enspace .$ & $0.6746802\,\mathrm n$ & \multicolumn{5}{@{\quad}c@{}}{ferner}\\
\cline{1-2}
$\log\frac{r}{R}\enspace .\ .\ .\ .\ .\ .\ .$ & $0.4394566$ & \multicolumn{5}{@{\quad}l@{}}{$\log r=\log R+\log\dfrac{r}{R}=0.4320724$}\\
\end{tabular*}
\end{minipage}\par
\par\clearpage
\noindent\SourceLeaf{96}%
\parbox[t]{\dimexpr\linewidth-12mm\relax}{\centering\fontsize{6.2}{8.8}\selectfont\bfseries EINES HIMMELSKÖRPERS DEM ORTE SEINES KNOTENS DER NEIGUNG DER BAHN ETC.\par}%
\makebox[12mm][r]{{\large\bfseries 91}}\par\nobreak\smallskip
\par\medskip\noindent\begin{minipage}[t]{\textwidth}
{\centering $6^{0}$.\par}
\smallskip\fontsize{8.7}{10.5}\selectfont
\noindent\begin{tabular*}{\linewidth}{@{}l@{\enspace}l@{\extracolsep{\fill}\quad}l@{\extracolsep{0pt}\enspace}r@{\enspace}r@{\enspace}r@{\enspace}r@{}}
$\log\operatorname{tang}i\enspace .\ .\ .\ .\ .$ & $9.2729872$ & & & & &\\
$\log\sin(\alpha-\AscendingNode)\enspace .\ .\ .\ .$ & $9.6658973\,\mathrm n$ & & & & &\\
\cline{1-2}
$\log\operatorname{tang}F\enspace .\ .\ .\ .\ .$ & $8.9388845\,\mathrm n$ & daher & $F=-$ & $4^{0}$ & $57'$ & $53''955$\\
\cline{1-2}
$\log\cos F\enspace .\ .\ .\ .\ .$ & $9.9983674$ &  & $F-\ObservedBeta=-$ & $1$ & $51$ & $20.394$\\
$\log\sin\ObservedBeta\enspace .\ .\ .\ .\ .\ .$ & $8.7343300\,\mathrm n$ & & & & &\\
$\log\sin(\odot-\AscendingNode)\enspace .\ .\ .\ .$ & $9.5348776\,\mathrm n$ & & & & &\\
$\text{Compl.}\ \log\sin(F-\ObservedBeta)\enspace .$ & $1.4896990\,\mathrm n$ & & & & &\\
$\text{Compl.}\ \log\sin(\upsilon-\AscendingNode)\enspace .$ & $0.6746802\,\mathrm n$ & & & & &\\
$\text{Compl.}\ \log\cos i\enspace .\ .\ .$ & $0.0075025$ & & & & &\\
\cline{1-2}
$\log\frac{r}{R}\enspace .\ .\ .\ .\ .\ .\ .$ & $0.4394567$ & \multicolumn{5}{@{\quad}l@{}}{nahe wie vorher.}\\
\end{tabular*}
\end{minipage}\par
\par\medskip\noindent\begin{minipage}[t]{\textwidth}
{\centering $7^{0}$.\par}
\smallskip\fontsize{8.7}{10.5}\selectfont
\noindent\begin{tabular*}{\linewidth}{@{}l@{\enspace}l@{\extracolsep{\fill}\quad}l@{\extracolsep{0pt}\enspace}r@{\enspace}r@{\enspace}r@{\enspace}r@{}}
$\log\cos i\enspace .\ .\ .\ .\ .\ .$ & $9.9924975$ & & & & &\\
$\log\operatorname{tang}(\upsilon-\AscendingNode)\enspace .\ .\ .$ & $9.3352577\,\mathrm n$ & & & & &\\
\cline{1-2}
$\log\operatorname{tang}G\enspace .\ .\ .\ .\ .$ & $9.3277552\,\mathrm n$ & also & $G=-$ & $12^{0}$ & $0'$ & $27''118$\\
\cline{1-2}
$\log\cos G\enspace .\ .\ ,\ .\ .$ & $9.9903922$ &  & $\alpha-\AscendingNode-G=-$ & $15$ & $35$ & $42.492$\\
$\log\sin(\odot-\alpha)\enspace .\ .\ .\ .$ & $9.8686173\,\mathrm n$ & & & & &\\
$\text{Compl.}\ \log\sin(\alpha-\AscendingNode-G)$ & $0.5705092\,\mathrm n$ & & & & &\\
$\text{Compl.}\ \log\cos(\upsilon-\AscendingNode)\enspace .$ & $0.0099379$ & & & & &\\
\cline{1-2}
$\log\frac{r}{R}\enspace .\ .\ .\ .\ .\ .\ .$ & $0.4394566$ & \multicolumn{5}{@{\quad}l@{}}{wie oben.}\\
\end{tabular*}
\end{minipage}\par
\par\medskip\noindent\begin{minipage}[t]{\textwidth}
{\centering $8^{0}$.\par}
\smallskip\fontsize{8.7}{10.5}\selectfont
\noindent\begin{tabular*}{\linewidth}{@{}l@{\enspace}l@{\extracolsep{\fill}\quad}l@{\extracolsep{0pt}\enspace}r@{\enspace}r@{\enspace}r@{\enspace}r@{}}
$\log\operatorname{tang}(\alpha-\AscendingNode)\enspace .\ .\ .$ & $9.7183744\,\mathrm n$ & & & & &\\
$\log\cos i\enspace .\ .\ .\ .\ .\ .$ & $9.9924975$ & & & & &\\
\cline{1-2}
$\log\operatorname{tang}H\enspace .\ .\ .\ .\ .$ & $9.7258769\,\mathrm n$ & folglich & $H=-$ & $28^{0}$ & $0'$ & $39''879$\\
\cline{1-2}
$\log\sin H\enspace .\ .\ .\ .\ .$ & $9.6717672\,\mathrm n$ &  & $H-(\upsilon-\AscendingNode)=-$ & $15$ & $48$ & $1.937$\\
$\log\sin(\odot-\alpha)\enspace .\ .\ .\ .$ & $9.8686173\,\mathrm n$ & & & & &\\
$\text{Compl.}\ \log\sin(H-(\upsilon-\AscendingNode))$ & $0.5649695\,\mathrm n$ & & & & &\\
$\text{Compl.}\ \log\sin(\alpha-\AscendingNode)\enspace .$ & $0.3341027\,\mathrm n$ & & & & &\\
\cline{1-2}
$\log\frac{r}{R}\enspace .\ .\ .\ .\ .\ .\ .$ & $0.4394567$ & \multicolumn{5}{@{\quad}l@{}}{wie vorher.}\\
\end{tabular*}
\end{minipage}\par
\par\GaussVIGathering{}{12\textsuperscript{*}}
\par\clearpage
\noindent\SourceLeaf{97}%
\makebox[12mm][l]{{\large\bfseries 92}}%
\parbox[t]{\dimexpr\linewidth-24mm\relax}{\centering\fontsize{7.05}{9.2}\selectfont\bfseries VORSCHRIFTEN UM AUS DER GEOCENTRISCHEN LÄNGE UND BREITE\par}%
\makebox[12mm][r]{}\par\nobreak\smallskip
\par\medskip\noindent\begin{minipage}[t]{\textwidth}
{\centering $9^{0}$.\par}
\smallskip\fontsize{8.7}{10.5}\selectfont
\noindent\begin{tabular*}{\linewidth}{@{}l@{\enspace}l@{\extracolsep{\fill}\quad}l@{\extracolsep{0pt}\enspace}r@{\enspace}r@{\enspace}r@{\enspace}r@{}}
$\log\operatorname{tang}\ObservedBeta\enspace .\ .\ .\ .\ .$ & $8.7349698\,\mathrm n$ & & & & &\\
$\text{Compl.}\ \log\sin i\enspace .\ .\ .$ & $0.7345153$ & & & & &\\
$\text{Compl.}\ \log\cos(\alpha-\AscendingNode)\enspace .$ & $0.0524771$ & & & & &\\
\cline{1-2}
$\log\operatorname{tang}I\enspace .\ .\ .\ .\ .$ & $9.5219622\,\mathrm n$ & hieraus & $I=-$ & $18^{0}$ & $23'$ & $55''334$\\
\cline{1-2}
$\log\sin I\enspace .\ .\ .\ .\ .\ .$ & $9.4991749\,\mathrm n$ &  & $\upsilon-\AscendingNode-I=$ & $6$ & $11$ & $17.392$\\
$\log\cos(\odot-\AscendingNode)\enspace .\ .\ .$ & $9.9728762\,\mathrm n$ & & & & &\\
$\text{Compl.}\ \log\sin(\upsilon-\AscendingNode-I)$ & $0.9674054$ & & & & &\\
\cline{1-2}
$\log\frac{r}{R}\enspace .\ .\ .\ .\ .\ .\ .$ & $0.4394565$ & \multicolumn{5}{@{\quad}l@{}}{wie vorhin.}\\
\end{tabular*}
\end{minipage}\par
\par\medskip\noindent\begin{minipage}[t]{\textwidth}
{\centering $10^{0}$.\par}
{\centering In der Nähe des Knotens weniger scharf.\par}
\smallskip\fontsize{8.7}{10.5}\selectfont
\noindent\begin{tabular*}{\linewidth}{@{}l@{\enspace}l@{\extracolsep{\fill}\quad}l@{\extracolsep{0pt}\enspace}r@{\enspace}r@{\enspace}r@{\enspace}r@{}}
$\log\sin i\enspace .\ .\ .\ .\ .\ .$ & $9.2654847$ & & & & &\\
$\log\cos(\alpha-\AscendingNode)\enspace .\ .\ .\ .$ & $9.9475229$ & & & & &\\
$\log\operatorname{tang}(\upsilon-\AscendingNode)\enspace .\ .\ .\ .$ & $9.3352577\,\mathrm n$ & & & & &\\
\cline{1-2}
$\log\operatorname{tang}K\enspace .\ .\ .\ .\ .$ & $8.5482653\,\mathrm n$ & also & $K=-$ & $2^{0}$ & $1'$ & $26''344$\\
\cline{1-2}
$\log\cos K\enspace .\ .\ .\ .\ .$ & $9.9997290$ &  & $K-\ObservedBeta=$ & $1$ & $5$ & $7.217$\\
$\log\sin\ObservedBeta\enspace .\ .\ .\ .\ .\ .$ & $8.7343300\,\mathrm n$ & & & & &\\
$\log\cos(\odot-\AscendingNode)\enspace .\ .\ .$ & $9.9728762\,\mathrm n$ & & & & &\\
$\text{Compl.}\ \log\sin(K-\ObservedBeta)\enspace .$ & $1.7225836$ & & & & &\\
$\text{Compl.}\ \log\cos(\upsilon-\AscendingNode)\enspace .$ & $0.0099379$ & & & & &\\
\cline{1-2}
$\log\frac{r}{R}\enspace .\ .\ .\ .\ .\ .\ .$ & $0.4394567$ & \multicolumn{5}{@{\quad}l@{}}{wie vorhin.}\\
\end{tabular*}
\end{minipage}\par
\par\medskip\noindent\begin{minipage}[t]{\textwidth}
{\centering $11^{0}$.\par}
\[C+\odot-\alpha=219^{0}\ 59'\ 25''474\]
\smallskip\fontsize{8.7}{10.5}\selectfont
\noindent\begin{tabular*}{\linewidth}{@{}l@{\enspace}l@{\extracolsep{\fill}\quad}l@{\extracolsep{0pt}\enspace}r@{\enspace}r@{\enspace}r@{\enspace}r@{}}
$\log\sin C\enspace .\ .\ .\ .\ .$ & $9.1243583\,\mathrm n$ & & & & &\\
$\log\sin(\odot-\alpha)\enspace .\ .\ .\ .$ & $9.8686173\,\mathrm n$ & & & & &\\
$\text{Compl.}\ \log\cos(C+\odot-\alpha)$ & $0.1156850\,\mathrm n$ & & & & &\\
$\text{Compl.}\ \log\operatorname{tang}(\odot-\AscendingNode)$ & $0.4379986$ & & & & &\\
$\text{Compl.}\ \log\cos i\enspace .\ .\ .$ & $0.0075025$ & & & & &\\
\cline{1-2}
$\log\operatorname{tang}L\enspace .\ .\ .\ .\ .$ & $9.5541617\,\mathrm n$ & also & $L=-$ & $19^{0}$ & $42'$ & $32''533$\\
\cline{1-2}
$\log\sin L\enspace .\ .\ .\ .\ .\ .$ & $9.5279439\,\mathrm n$ &  & $\upsilon-\AscendingNode-L=$ & $7$ & $29$ & $54.591$\\
$\text{Compl.}\ \log\sin(\upsilon-\AscendingNode-L)$ & $0.8843888$ & & & & &\\
$\text{Compl.}\ \log\cos(\odot-\AscendingNode)\enspace .$ & $0.0271238\,\mathrm n$ & & & & &\\
\cline{1-2}
$\log\frac{r}{R}\enspace .\ .\ .\ .\ .\ .\ .$ & $0.4394565$ & \multicolumn{5}{@{\quad}l@{}}{wie zuvor.}\\
\end{tabular*}
\end{minipage}\par
\par\clearpage
\noindent\SourceLeaf{98}%
\parbox[t]{\dimexpr\linewidth-12mm\relax}{\centering\fontsize{6.2}{8.8}\selectfont\bfseries EINES HIMMELSKÖRPERS DEM ORTE SEINES KNOTENS DER NEIGUNG DER BAHN ETC.\par}%
\makebox[12mm][r]{{\large\bfseries 93}}\par\nobreak\smallskip
\par\medskip\noindent\begin{minipage}[t]{\textwidth}
{\centering $12^{0}$.\par}
\[D+\odot-\AscendingNode=178^{0}\ 2'\ 31''738\]
\smallskip\fontsize{8.7}{10.5}\selectfont
\noindent\begin{tabular*}{\linewidth}{@{}l@{\enspace}l@{\extracolsep{\fill}\quad}l@{\extracolsep{0pt}\enspace}r@{\enspace}r@{\enspace}r@{\enspace}r@{}}
$\log\sin D\enspace .\ .\ .\ .\ .$ & $9.5735295\,\mathrm n$ & & & & &\\
$\log\cos(\odot-\AscendingNode)\enspace .\ .\ .$ & $9.9728762\,\mathrm n$ & & & & &\\
$\text{Compl.}\ \log\cos(D+\odot-\AscendingNode)$ & $0.0002536\,\mathrm n$ & & & & &\\
$\text{Compl.}\ \log\cos i\enspace .\ .\ .$ & $0.0075025$ & & & & &\\
\cline{1-2}
$\log\operatorname{tang}(M=L)\enspace .\ .\ .$ & $9.5541618\,\mathrm n$ & \multicolumn{5}{@{\quad}l@{}}{wie oben in $11^{0}$.}\\
\end{tabular*}
\par\smallskip\hspace*{1.2em}Der übrige Theil der Rechnung eben so wie dort.
\end{minipage}\par
\par\medskip\noindent\begin{minipage}[t]{\textwidth}
{\centering $13^{0}$.\par}
\smallskip\fontsize{8.7}{10.5}\selectfont
\noindent\begin{tabular}{@{}l@{\enspace}l@{}}
$\log r\enspace .\ .\ .\ .\ .\ .\ .$ & $0.4320724$\\
$\log\sin(\upsilon-\AscendingNode)\enspace .\ .\ .\ .$ & $9.3253198\,\mathrm n$\\
$\log\sin i\enspace .\ .\ .\ .\ .\ .$ & $9.2654847$\\
$\text{Compl.}\ \log\sin\ObservedBeta\enspace .\ .\ .$ & $1.2656700\,\mathrm n$\\
\hline
$\log\Delta\enspace .\ .\ .\ .\ .\ .\ .$ & $0.2885469$\\
\end{tabular}
\end{minipage}\par
\par\bigskip\bigskip
\centerline{\rule{.25\textwidth}{.3pt}}
\clearpage
\endgroup
